\documentclass[12pt]{article}

%----------------- Packages ------------------
\usepackage[utf8]{inputenc}
\usepackage[T1]{fontenc}
\usepackage{amsmath, amssymb}
\usepackage{geometry}
\usepackage{xcolor}
\usepackage{titlesec}
\usepackage{fancyhdr}
\usepackage[breakable]{tcolorbox}
\usepackage{graphicx}
\usepackage{tikz}
\usepackage{float}
\usetikzlibrary{arrows.meta, decorations.markings}

%----------------- Page Setup -----------------
\geometry{a4paper, margin=2.5cm}
\pagestyle{fancy}
\fancyhf{}
\rhead{Final Exam — \textbf{2015}}
\lhead{Rigid Body Mechanics}
\cfoot{\thepage}

%----------------- Title Styling --------------
\titleformat{\section}{\normalfont\Large\bfseries}{Exercise \thesection:}{1em}{}
\titleformat{\subsection}{\normalfont\bfseries}{Answer:}{1em}{}

%----------------- Custom Boxes ----------------
\tcbuselibrary{listingsutf8}
\newtcolorbox{answerbox}{
  colback=gray!10,
  colframe=black,
  fonttitle=\bfseries,
  title=Answer Area,
  breakable,
  before skip=10pt,
  after skip=10pt
}

%----------------- Document Start --------------
\begin{document}

%----------------- Exam Info -------------------
\begin{center}
  \Large\textbf{UNIVERSITY NAME} \\[1em]
  \large\textit{Rigid Body Mechanics — Final Exam} \\[0.5em]
  \large\textit{Year: 2015} \\[2em]
\end{center}

\vspace{0.5cm}

%----------------- EXERCISE 1 ------------------
\section{}
\begin{enumerate}
    \item Show that the velocity vector field of a rigid body defines a torsor.
    
    \item Show that the angular momentum is an antisymmetric vector field.
    
    \item Show the following formula: for any point \(A\) belonging to the rigid body \(S\) of mass \(m\) and center of mass \(G\),
    \[
    L_A(S/R_0) = \Pi_A(S) \cdot \hat{\Omega}(S/R_0) + m \cdot \overrightarrow{AG} \times \overrightarrow{V}(A/R_0).
    \]
    
    \item 
    \begin{enumerate}
        \item Write the theorem of angular momentum about a point \(A\) that is mobile in \(R_0\).
        \item Prove this theorem.
    \end{enumerate}
\end{enumerate}

\begin{answerbox}
\begin{enumerate}
    \item 
    \item 
    \item 
    \item 
    \begin{enumerate}
        \item 
        \item 
    \end{enumerate}
\end{enumerate}
\end{answerbox}

\newpage
%----------------- EXERCISE 2 ------------------
\section{}
A physical pendulum \(S\), consisting of a homogeneous rod of length \(B_a\) and negligible mass, welded to a homogeneous ball of mass \(m\), radius \(a\), and center of mass \(G\), is in motion relative to a fixed Galilean reference frame \(R_0(X_0, Y_0, Z_0)\) with basis \((\bar{x}_0, \bar{y}_0, \bar{z}_0)\).

The system is moving in the vertical plane \(X_0OY_0\) (with \(G \in X_0OY_0\)) and \(A\) is a mobile point. The motion is frictionless at \(A\).  
\(R\) denotes the reaction force exerted on the system at point \(A\).

The position of point \(A\) on the axis \(OY_0\) is given by the coordinate \(y\) (with \(\overrightarrow{OA} = y \overrightarrow{y_0}\)) and the position of the system relative to the vertical is defined by the angle \(\theta\). The system is displaced from its equilibrium position and then released without initial velocity.

\begin{enumerate}
    \item Calculate the inertia matrix of the ball at point \(G\).
    
    \item Determine the velocity vector of point \(G\) in \(R_0\).
    
    \item By applying the dynamic resultant theorem, show that \(y\) and \(\theta\) are related by the following relation: \(y = -6a \dot{\theta} \cos\theta\). Deduce that the kinetic resultant of the system is given by:
    \[
    P(S/R_0) = -6m a \dot{\theta} \sin\theta \overrightarrow{x_0}
    \]
    
    \item By applying the kinetic energy theorem to the solid, determine the first integral of energy and deduce the differential equation of motion of the solid.
    
    \item Retrieve the differential equation of motion of the solid by applying the angular momentum theorem to the solid, at point \(A\), in the reference frame \(R_0\).
    
    \item Determine the period of small oscillations (keeping only first-order terms in \(\theta\) and \(\dot{\theta}\)).
\end{enumerate}

\newpage

\begin{answerbox}
\begin{enumerate}
  \item 
  \item 
  \item 
  \item 
  \item 
  \item 
\end{enumerate}
\end{answerbox}

%----------------- END ------------------
\end{document}